Concentrated liquidity provision on Uniswap~v3 turns a passive role that only
earns fees into an active control problem. A liquidity provider (LP) must
decide when to recenter its price band, how wide to keep it, and how much of
the position's price exposure to hedge, while fee income competes with
impermanent loss, loss-versus-rebalancing (LVR), gas and slippage. Empirical
studies show that fees frequently fail to cover these losses. Whether one
learned policy can manage all three decisions jointly under realistic
frictions, and where active management stops paying, remains open. We present
\emph{Undertow}, a research testbed with two engines: a replay backtester whose
fee accounting stays in exact integer arithmetic across two years of on-chain
Uniswap~v3 events, and a fast Gymnasium simulator. We formulate liquidity
provision as a risk-sensitive Markov decision process (MDP) with impulse
control over a joint action space of timing, width and hedge ratio. Four claims structure
the study: a negative limit that locates where fees fall short of LVR and
rebalancing costs; a methods contribution documenting why naive PPO degenerates
on this MDP and which fixes (band geometry in the observation, an
excess-versus-HODL reward, environment diversity, and execution at decision
time) are necessary; evidence from windows held out from training on whether a
single risk-sensitive policy dominates the best static width and a hedged
HODL benchmark matched in beta; and a map of the active/passive boundary in
terms of volatility, fee tier, gas and capital.
\pending{The abstract will close with the quantitative headline and the C3/C4
verdicts once the locked experiments across seeds land.}
